% =====================================================================
%  Anatomía de un LLM — Fase 0 — Módulo 0.4
%  Trigonometría y números complejos
%  Compilar: latexmk -r ../latexmkrc P04-trigonometria-complejos.tex
% =====================================================================
\documentclass[11pt,a4paper,twoside,openany]{book}
\usepackage{../estilo-apuntes}

\newcommand{\ApuntesCurso}{Anatomía de un LLM}
\newcommand{\ApuntesModulo}{Módulo 0.4}

\newcommand{\T}{^{\top}}
\newcommand{\C}{\mathbb{C}}
\newcommand{\Z}{\mathbb{Z}}
\newcommand{\N}{\mathbb{N}}
\newcommand{\Real}{\operatorname{Re}}
\newcommand{\Imag}{\operatorname{Im}}
\newcommand{\Arg}{\operatorname{Arg}}
\newcommand{\PE}{\operatorname{PE}}
% Anclas de hyperref únicas por tema (evita destinos duplicados en el PDF)
\makeatletter
\@for\p@env:=definicion,teorema,proposicion,lema,corolario,ejemplo,observacion\do{%
  \expandafter\def\csname theH\p@env\endcsname{\thechapter.\arabic{definicion}}}
\makeatother

\hypersetup{pdftitle={Módulo 0.4. Trigonometría y números complejos (resumen)},pdfauthor={Anatomía de un LLM}}

\makeatletter
\newlabel{ej:arco}{{1.3}{1}{}{ejemplo.1.3}{}}
\newlabel{ej:arco@cref}{{[ejemplo][3][]1.3}{[1][1][]1}}
\newlabel{fig:circunferencia}{{1.1}{3}{La circunferencia unidad. El arco de longitud $\theta $ desde $(1,0)$ llega al punto $P(\theta )$; sus coordenadas son $\cos \theta $ y $\sin \theta $. Los números romanos indican los cuadrantes}{figure.1.1}{}}
\newlabel{fig:circunferencia@cref}{{[figure][1][1]1.1}{[1][2][]3}}
\newlabel{ej:reduccion}{{1.12}{4}{Valores por reducción}{ejemplo.1.12}{}}
\newlabel{ej:reduccion@cref}{{[ejemplo][12][]1.12}{[1][3][]4}}
\newlabel{eje:1-conversion}{{1.1}{4}{}{ejercicio.1.1}{}}
\newlabel{eje:1-conversion@cref}{{[ejercicio][1][1]1.1}{[1][4][]4}}
\newlabel{eje:1-cuadrantes}{{1.2}{4}{}{ejercicio.1.2}{}}
\newlabel{eje:1-cuadrantes@cref}{{[ejercicio][2][1]1.2}{[1][4][]4}}
\newlabel{eje:1-vf}{{1.3}{4}{}{ejercicio.1.3}{}}
\newlabel{eje:1-vf@cref}{{[ejercicio][3][1]1.3}{[1][4][]4}}
\newlabel{eje:1-arco}{{1.4}{4}{}{ejercicio.1.4}{}}
\newlabel{eje:1-arco@cref}{{[ejercicio][4][1]1.4}{[1][4][]4}}
\newlabel{eje:1-exactos}{{1.5}{5}{}{ejercicio.1.5}{}}
\newlabel{eje:1-exactos@cref}{{[ejercicio][5][1]1.5}{[1][4][]5}}
\newlabel{eje:1-seis-razones}{{1.6}{5}{}{ejercicio.1.6}{}}
\newlabel{eje:1-seis-razones@cref}{{[ejercicio][6][1]1.6}{[1][5][]5}}
\newlabel{eje:1-elevacion}{{1.7}{5}{}{ejercicio.1.7}{}}
\newlabel{eje:1-elevacion@cref}{{[ejercicio][7][1]1.7}{[1][5][]5}}
\newlabel{eje:1-coterminales}{{1.8}{5}{}{ejercicio.1.8}{}}
\newlabel{eje:1-coterminales@cref}{{[ejercicio][8][1]1.8}{[1][5][]5}}
\newlabel{eje:1-rope}{{1.9}{5}{}{ejercicio.1.9}{}}
\newlabel{eje:1-rope@cref}{{[ejercicio][9][1]1.9}{[1][5][]5}}
\newlabel{eje:1-simplificar}{{1.10}{5}{}{ejercicio.1.10}{}}
\newlabel{eje:1-simplificar@cref}{{[ejercicio][10][1]1.10}{[1][5][]5}}
\newlabel{eje:1-sinx-x-tanx}{{1.11}{5}{}{ejercicio.1.11}{}}
\newlabel{eje:1-sinx-x-tanx@cref}{{[ejercicio][11][1]1.11}{[1][5][]5}}
\newlabel{eje:1-abs-sin}{{1.12}{5}{}{ejercicio.1.12}{}}
\newlabel{eje:1-abs-sin@cref}{{[ejercicio][12][1]1.12}{[1][5][]5}}
\newlabel{eje:1-sin-mas-cos}{{1.13}{5}{}{ejercicio.1.13}{}}
\newlabel{eje:1-sin-mas-cos@cref}{{[ejercicio][13][1]1.13}{[1][5][]5}}
\newlabel{eje:1-periodo-tan}{{1.14}{5}{}{ejercicio.1.14}{}}
\newlabel{eje:1-periodo-tan@cref}{{[ejercicio][14][1]1.14}{[1][5][]5}}
\newlabel{eje:1-poligono}{{1.15}{5}{}{ejercicio.1.15}{}}
\newlabel{eje:1-poligono@cref}{{[ejercicio][15][1]1.15}{[1][5][]5}}
\newlabel{ej:verificar}{{2.2}{6}{}{ejemplo.2.2}{}}
\newlabel{ej:verificar@cref}{{[ejemplo][2][]2.2}{[1][6][]6}}
\newlabel{ej:cos15}{{2.5}{7}{}{ejemplo.2.5}{}}
\newlabel{ej:cos15@cref}{{[ejemplo][5][]2.5}{[1][7][]7}}
\newlabel{ej:mitad}{{2.9}{8}{}{ejemplo.2.9}{}}
\newlabel{ej:mitad@cref}{{[ejemplo][9][]2.9}{[1][8][]8}}
\newlabel{eje:2-vf}{{2.1}{9}{}{ejercicio.2.1}{}}
\newlabel{eje:2-vf@cref}{{[ejercicio][1][2]2.1}{[1][8][]9}}
\newlabel{eje:2-signo-mitad}{{2.2}{9}{}{ejercicio.2.2}{}}
\newlabel{eje:2-signo-mitad@cref}{{[ejercicio][2][2]2.2}{[1][9][]9}}
\newlabel{eje:2-exactos}{{2.3}{9}{}{ejercicio.2.3}{}}
\newlabel{eje:2-exactos@cref}{{[ejercicio][3][2]2.3}{[1][9][]9}}
\newlabel{eje:2-cuadrantes}{{2.4}{9}{}{ejercicio.2.4}{}}
\newlabel{eje:2-cuadrantes@cref}{{[ejercicio][4][2]2.4}{[1][9][]9}}
\newlabel{eje:2-doble}{{2.5}{9}{}{ejercicio.2.5}{}}
\newlabel{eje:2-doble@cref}{{[ejercicio][5][2]2.5}{[1][9][]9}}
\newlabel{eje:2-pi8}{{2.6}{9}{}{ejercicio.2.6}{}}
\newlabel{eje:2-pi8@cref}{{[ejercicio][6][2]2.6}{[1][9][]9}}
\newlabel{eje:2-producto}{{2.7}{9}{}{ejercicio.2.7}{}}
\newlabel{eje:2-producto@cref}{{[ejercicio][7][2]2.7}{[1][9][]9}}
\newlabel{eje:2-potencias}{{2.8}{9}{}{ejercicio.2.8}{}}
\newlabel{eje:2-potencias@cref}{{[ejercicio][8][2]2.8}{[1][9][]9}}
\newlabel{eje:2-verificar}{{2.9}{9}{}{ejercicio.2.9}{}}
\newlabel{eje:2-verificar@cref}{{[ejercicio][9][2]2.9}{[1][9][]9}}
\newlabel{eje:2-triple}{{2.10}{9}{}{ejercicio.2.10}{}}
\newlabel{eje:2-triple@cref}{{[ejercicio][10][2]2.10}{[1][9][]9}}
\newlabel{eje:2-cos36}{{2.11}{9}{}{ejercicio.2.11}{}}
\newlabel{eje:2-cos36@cref}{{[ejercicio][11][2]2.11}{[1][9][]9}}
\newlabel{eje:2-suma-senos}{{2.12}{9}{}{ejercicio.2.12}{}}
\newlabel{eje:2-suma-senos@cref}{{[ejercicio][12][2]2.12}{[1][9][]9}}
\newlabel{eje:2-tan-triangulo}{{2.13}{10}{}{ejercicio.2.13}{}}
\newlabel{eje:2-tan-triangulo@cref}{{[ejercicio][13][2]2.13}{[1][10][]10}}
\newlabel{eje:2-chebyshev}{{2.14}{10}{}{ejercicio.2.14}{}}
\newlabel{eje:2-chebyshev@cref}{{[ejercicio][14][2]2.14}{[1][10][]10}}
\newlabel{eje:2-potencia-giro}{{2.15}{10}{}{ejercicio.2.15}{}}
\newlabel{eje:2-potencia-giro@cref}{{[ejercicio][15][2]2.15}{[1][10][]10}}
\newlabel{ej:lal}{{3.6}{12}{Caso LAL}{ejemplo.3.6}{}}
\newlabel{ej:lal@cref}{{[ejemplo][6][]3.6}{[1][12][]12}}
\newlabel{ej:angulo}{{3.10}{13}{}{ejemplo.3.10}{}}
\newlabel{ej:angulo@cref}{{[ejemplo][10][]3.10}{[1][13][]13}}
\newlabel{eje:3-que-teorema}{{3.1}{14}{}{ejercicio.3.1}{}}
\newlabel{eje:3-que-teorema@cref}{{[ejercicio][1][3]3.1}{[1][14][]14}}
\newlabel{eje:3-imposible}{{3.2}{14}{}{ejercicio.3.2}{}}
\newlabel{eje:3-imposible@cref}{{[ejercicio][2][3]3.2}{[1][14][]14}}
\newlabel{eje:3-obtuso}{{3.3}{14}{}{ejercicio.3.3}{}}
\newlabel{eje:3-obtuso@cref}{{[ejercicio][3][3]3.3}{[1][14][]14}}
\newlabel{eje:3-ala}{{3.4}{14}{}{ejercicio.3.4}{}}
\newlabel{eje:3-ala@cref}{{[ejercicio][4][3]3.4}{[1][14][]14}}
\newlabel{eje:3-ambiguo}{{3.5}{14}{}{ejercicio.3.5}{}}
\newlabel{eje:3-ambiguo@cref}{{[ejercicio][5][3]3.5}{[1][14][]14}}
\newlabel{eje:3-lal}{{3.6}{14}{}{ejercicio.3.6}{}}
\newlabel{eje:3-lal@cref}{{[ejercicio][6][3]3.6}{[1][14][]14}}
\newlabel{eje:3-barcos}{{3.7}{14}{}{ejercicio.3.7}{}}
\newlabel{eje:3-barcos@cref}{{[ejercicio][7][3]3.7}{[1][14][]14}}
\newlabel{eje:3-lll}{{3.8}{14}{}{ejercicio.3.8}{}}
\newlabel{eje:3-lll@cref}{{[ejercicio][8][3]3.8}{[1][14][]14}}
\newlabel{eje:3-angulo-vectores}{{3.9}{14}{}{ejercicio.3.9}{}}
\newlabel{eje:3-angulo-vectores@cref}{{[ejercicio][9][3]3.9}{[1][14][]14}}
\newlabel{eje:3-coseno-normas}{{3.10}{14}{}{ejercicio.3.10}{}}
\newlabel{eje:3-coseno-normas@cref}{{[ejercicio][10][3]3.10}{[1][14][]14}}
\newlabel{eje:3-similitud}{{3.11}{14}{}{ejercicio.3.11}{}}
\newlabel{eje:3-similitud@cref}{{[ejercicio][11][3]3.11}{[1][14][]14}}
\newlabel{eje:3-circunradio}{{3.12}{15}{}{ejercicio.3.12}{}}
\newlabel{eje:3-circunradio@cref}{{[ejercicio][12][3]3.12}{[1][14][]15}}
\newlabel{eje:3-heron}{{3.13}{15}{}{ejercicio.3.13}{}}
\newlabel{eje:3-heron@cref}{{[ejercicio][13][3]3.13}{[1][15][]15}}
\newlabel{eje:3-proyecciones}{{3.14}{15}{}{ejercicio.3.14}{}}
\newlabel{eje:3-proyecciones@cref}{{[ejercicio][14][3]3.14}{[1][15][]15}}
\newlabel{eje:3-mayor-lado}{{3.15}{15}{}{ejercicio.3.15}{}}
\newlabel{eje:3-mayor-lado@cref}{{[ejercicio][15][3]3.15}{[1][15][]15}}
\newlabel{fig:graficas}{{4.1}{17}{Gráficas de $\sin x$ (trazo continuo) y $\cos x$ (discontinuo) en $[-\pi ,2\pi ]$}{figure.4.1}{}}
\newlabel{fig:graficas@cref}{{[figure][1][4]4.1}{[1][16][]17}}
\newlabel{ej:inversas}{{4.7}{18}{}{ejemplo.4.7}{}}
\newlabel{ej:inversas@cref}{{[ejemplo][7][]4.7}{[1][18][]18}}
\newlabel{ej:cuadratica}{{4.9}{18}{Ecuación de segundo grado en el seno}{ejemplo.4.9}{}}
\newlabel{ej:cuadratica@cref}{{[ejemplo][9][]4.9}{[1][18][]18}}
\newlabel{eje:4-parametros}{{4.1}{19}{}{ejercicio.4.1}{}}
\newlabel{eje:4-parametros@cref}{{[ejercicio][1][4]4.1}{[1][19][]19}}
\newlabel{eje:4-dominios}{{4.2}{19}{}{ejercicio.4.2}{}}
\newlabel{eje:4-dominios@cref}{{[ejercicio][2][4]4.2}{[1][19][]19}}
\newlabel{eje:4-exactos}{{4.3}{19}{}{ejercicio.4.3}{}}
\newlabel{eje:4-exactos@cref}{{[ejercicio][3][4]4.3}{[1][19][]19}}
\newlabel{eje:4-basicas}{{4.4}{19}{}{ejercicio.4.4}{}}
\newlabel{eje:4-basicas@cref}{{[ejercicio][4][4]4.4}{[1][19][]19}}
\newlabel{eje:4-cuadratica}{{4.5}{19}{}{ejercicio.4.5}{}}
\newlabel{eje:4-cuadratica@cref}{{[ejercicio][5][4]4.5}{[1][19][]19}}
\newlabel{eje:4-doble}{{4.6}{19}{}{ejercicio.4.6}{}}
\newlabel{eje:4-doble@cref}{{[ejercicio][6][4]4.6}{[1][19][]19}}
\newlabel{eje:4-armonico}{{4.7}{19}{}{ejercicio.4.7}{}}
\newlabel{eje:4-armonico@cref}{{[ejercicio][7][4]4.7}{[1][19][]19}}
\newlabel{eje:4-cuadrado}{{4.8}{19}{}{ejercicio.4.8}{}}
\newlabel{eje:4-cuadrado@cref}{{[ejercicio][8][4]4.8}{[1][19][]19}}
\newlabel{eje:4-periodo-suma}{{4.9}{19}{}{ejercicio.4.9}{}}
\newlabel{eje:4-periodo-suma@cref}{{[ejercicio][9][4]4.9}{[1][19][]19}}
\newlabel{eje:4-pe}{{4.10}{19}{}{ejercicio.4.10}{}}
\newlabel{eje:4-pe@cref}{{[ejercicio][10][4]4.10}{[1][19][]19}}
\newlabel{eje:4-creciente}{{4.11}{20}{}{ejercicio.4.11}{}}
\newlabel{eje:4-creciente@cref}{{[ejercicio][11][4]4.11}{[1][19][]20}}
\newlabel{eje:4-cos-iguales}{{4.12}{20}{}{ejercicio.4.12}{}}
\newlabel{eje:4-cos-iguales@cref}{{[ejercicio][12][4]4.12}{[1][20][]20}}
\newlabel{eje:4-arctan-inverso}{{4.13}{20}{}{ejercicio.4.13}{}}
\newlabel{eje:4-arctan-inverso@cref}{{[ejercicio][13][4]4.13}{[1][20][]20}}
\newlabel{eje:4-machin}{{4.14}{20}{}{ejercicio.4.14}{}}
\newlabel{eje:4-machin@cref}{{[ejercicio][14][4]4.14}{[1][20][]20}}
\newlabel{eje:4-pe-lineal}{{4.15}{20}{}{ejercicio.4.15}{}}
\newlabel{eje:4-pe-lineal@cref}{{[ejercicio][15][4]4.15}{[1][20][]20}}
\newlabel{ej:potencias-i}{{5.3}{21}{Potencias de $i$}{ejemplo.5.3}{}}
\newlabel{ej:potencias-i@cref}{{[ejemplo][3][]5.3}{[1][21][]21}}
\newlabel{ej:division}{{5.6}{22}{}{ejemplo.5.6}{}}
\newlabel{ej:division@cref}{{[ejemplo][6][]5.6}{[1][22][]22}}
\newlabel{ej:raiz}{{5.10}{23}{}{ejemplo.5.10}{}}
\newlabel{ej:raiz@cref}{{[ejemplo][10][]5.10}{[1][23][]23}}
\newlabel{fig:argand}{{5.1}{24}{El plano complejo. El número $z=a+bi$ es el punto $(a,b)$; su módulo es la distancia al origen y su conjugado, la reflexión en el eje real}{figure.5.1}{}}
\newlabel{fig:argand@cref}{{[figure][1][5]5.1}{[1][23][]24}}
\newlabel{eje:5-partes}{{5.1}{24}{}{ejercicio.5.1}{}}
\newlabel{eje:5-partes@cref}{{[ejercicio][1][5]5.1}{[1][24][]24}}
\newlabel{eje:5-vf}{{5.2}{24}{}{ejercicio.5.2}{}}
\newlabel{eje:5-vf@cref}{{[ejercicio][2][5]5.2}{[1][24][]24}}
\newlabel{eje:5-potencias-i}{{5.3}{24}{}{ejercicio.5.3}{}}
\newlabel{eje:5-potencias-i@cref}{{[ejercicio][3][5]5.3}{[1][24][]24}}
\newlabel{eje:5-operaciones}{{5.4}{25}{}{ejercicio.5.4}{}}
\newlabel{eje:5-operaciones@cref}{{[ejercicio][4][5]5.4}{[1][24][]25}}
\newlabel{eje:5-inversos}{{5.5}{25}{}{ejercicio.5.5}{}}
\newlabel{eje:5-inversos@cref}{{[ejercicio][5][5]5.5}{[1][25][]25}}
\newlabel{eje:5-cuadraticas}{{5.6}{25}{}{ejercicio.5.6}{}}
\newlabel{eje:5-cuadraticas@cref}{{[ejercicio][6][5]5.6}{[1][25][]25}}
\newlabel{eje:5-suma-tripos}{{5.7}{25}{}{ejercicio.5.7}{}}
\newlabel{eje:5-suma-tripos@cref}{{[ejercicio][7][5]5.7}{[1][25][]25}}
\newlabel{eje:5-raices-cuadradas}{{5.8}{25}{}{ejercicio.5.8}{}}
\newlabel{eje:5-raices-cuadradas@cref}{{[ejercicio][8][5]5.8}{[1][25][]25}}
\newlabel{eje:5-cuadratica-compleja}{{5.9}{25}{}{ejercicio.5.9}{}}
\newlabel{eje:5-cuadratica-compleja@cref}{{[ejercicio][9][5]5.9}{[1][25][]25}}
\newlabel{eje:5-lugares}{{5.10}{25}{}{ejercicio.5.10}{}}
\newlabel{eje:5-lugares@cref}{{[ejercicio][10][5]5.10}{[1][25][]25}}
\newlabel{eje:5-escalar}{{5.11}{25}{}{ejercicio.5.11}{}}
\newlabel{eje:5-escalar@cref}{{[ejercicio][11][5]5.11}{[1][25][]25}}
\newlabel{eje:5-modulo-producto}{{5.12}{25}{}{ejercicio.5.12}{}}
\newlabel{eje:5-modulo-producto@cref}{{[ejercicio][12][5]5.12}{[1][25][]25}}
\newlabel{eje:5-igualdad-triangular}{{5.13}{25}{}{ejercicio.5.13}{}}
\newlabel{eje:5-igualdad-triangular@cref}{{[ejercicio][13][5]5.13}{[1][25][]25}}
\newlabel{eje:5-orden}{{5.14}{25}{}{ejercicio.5.14}{}}
\newlabel{eje:5-orden@cref}{{[ejercicio][14][5]5.14}{[1][25][]25}}
\newlabel{eje:5-blaschke}{{5.15}{25}{}{ejercicio.5.15}{}}
\newlabel{eje:5-blaschke@cref}{{[ejercicio][15][5]5.15}{[1][25][]25}}
\newlabel{ej:polar}{{6.3}{26}{}{ejemplo.6.3}{}}
\newlabel{ej:polar@cref}{{[ejemplo][3][]6.3}{[1][26][]26}}
\newlabel{ej:demoivre}{{6.10}{28}{Potencias y ángulo triple}{ejemplo.6.10}{}}
\newlabel{ej:demoivre@cref}{{[ejemplo][10][]6.10}{[1][28][]28}}
\newlabel{fig:raices}{{6.1}{29}{Las raíces sextas de la unidad, $\omega ^{k}=e^{ik\pi /3}$, son los vértices de un hexágono regular inscrito en la circunferencia unidad}{figure.6.1}{}}
\newlabel{fig:raices@cref}{{[figure][1][6]6.1}{[1][29][]29}}
\newlabel{eje:6-exponenciales}{{6.1}{30}{}{ejercicio.6.1}{}}
\newlabel{eje:6-exponenciales@cref}{{[ejercicio][1][6]6.1}{[1][30][]30}}
\newlabel{eje:6-vf}{{6.2}{30}{}{ejercicio.6.2}{}}
\newlabel{eje:6-vf@cref}{{[ejercicio][2][6]6.2}{[1][30][]30}}
\newlabel{eje:6-cuadrante}{{6.3}{30}{}{ejercicio.6.3}{}}
\newlabel{eje:6-cuadrante@cref}{{[ejercicio][3][6]6.3}{[1][30][]30}}
\newlabel{eje:6-polar}{{6.4}{30}{}{ejercicio.6.4}{}}
\newlabel{eje:6-polar@cref}{{[ejercicio][4][6]6.4}{[1][30][]30}}
\newlabel{eje:6-producto-polar}{{6.5}{30}{}{ejercicio.6.5}{}}
\newlabel{eje:6-producto-polar@cref}{{[ejercicio][5][6]6.5}{[1][30][]30}}
\newlabel{eje:6-potencias}{{6.6}{30}{}{ejercicio.6.6}{}}
\newlabel{eje:6-potencias@cref}{{[ejercicio][6][6]6.6}{[1][30][]30}}
\newlabel{eje:6-tripos}{{6.7}{30}{}{ejercicio.6.7}{}}
\newlabel{eje:6-tripos@cref}{{[ejercicio][7][6]6.7}{[1][30][]30}}
\newlabel{eje:6-raices}{{6.8}{31}{}{ejercicio.6.8}{}}
\newlabel{eje:6-raices@cref}{{[ejercicio][8][6]6.8}{[1][31][]31}}
\newlabel{eje:6-polinomios}{{6.9}{31}{}{ejercicio.6.9}{}}
\newlabel{eje:6-polinomios@cref}{{[ejercicio][9][6]6.9}{[1][31][]31}}
\newlabel{eje:6-giros}{{6.10}{31}{}{ejercicio.6.10}{}}
\newlabel{eje:6-giros@cref}{{[ejercicio][10][6]6.10}{[1][31][]31}}
\newlabel{eje:6-rope}{{6.11}{31}{}{ejercicio.6.11}{}}
\newlabel{eje:6-rope@cref}{{[ejercicio][11][6]6.11}{[1][31][]31}}
\newlabel{eje:6-suma-euler}{{6.12}{31}{}{ejercicio.6.12}{}}
\newlabel{eje:6-suma-euler@cref}{{[ejercicio][12][6]6.12}{[1][31][]31}}
\newlabel{eje:6-cos5}{{6.13}{31}{}{ejercicio.6.13}{}}
\newlabel{eje:6-cos5@cref}{{[ejercicio][13][6]6.13}{[1][31][]31}}
\newlabel{eje:6-potencias-raices}{{6.14}{31}{}{ejercicio.6.14}{}}
\newlabel{eje:6-potencias-raices@cref}{{[ejercicio][14][6]6.14}{[1][31][]31}}
\newlabel{eje:6-suma-cosenos}{{6.15}{31}{}{ejercicio.6.15}{}}
\newlabel{eje:6-suma-cosenos@cref}{{[ejercicio][15][6]6.15}{[1][31][]31}}
\newlabel{eje:6-autovalores}{{6.16}{32}{}{ejercicio.6.16}{}}
\newlabel{eje:6-autovalores@cref}{{[ejercicio][16][6]6.16}{[1][31][]32}}
\newlabel{ap:soluciones}{{A}{33}{Soluciones de los ejercicios}{appendix.A}{}}
\newlabel{ap:soluciones@cref}{{[appendix][1][2147483647]A}{[1][33][]33}}
\newlabel{sol:eje:1-conversion}{{A.1}{33}{Tema 1}{section*.4}{}}
\newlabel{sol:eje:1-conversion@cref}{{[subappendix][1][2147483647,1]A.1}{[1][33][]33}}
\newlabel{sol:eje:1-cuadrantes}{{A.1}{33}{Tema 1}{section*.5}{}}
\newlabel{sol:eje:1-cuadrantes@cref}{{[subappendix][1][2147483647,1]A.1}{[1][33][]33}}
\newlabel{sol:eje:1-vf}{{A.1}{33}{Tema 1}{section*.6}{}}
\newlabel{sol:eje:1-vf@cref}{{[subappendix][1][2147483647,1]A.1}{[1][33][]33}}
\newlabel{sol:eje:1-arco}{{A.1}{34}{Tema 1}{section*.7}{}}
\newlabel{sol:eje:1-arco@cref}{{[subappendix][1][2147483647,1]A.1}{[1][34][]34}}
\newlabel{sol:eje:1-exactos}{{A.1}{34}{Tema 1}{section*.8}{}}
\newlabel{sol:eje:1-exactos@cref}{{[subappendix][1][2147483647,1]A.1}{[1][34][]34}}
\newlabel{sol:eje:1-seis-razones}{{A.1}{34}{Tema 1}{section*.9}{}}
\newlabel{sol:eje:1-seis-razones@cref}{{[subappendix][1][2147483647,1]A.1}{[1][34][]34}}
\newlabel{sol:eje:1-elevacion}{{A.1}{35}{Tema 1}{section*.10}{}}
\newlabel{sol:eje:1-elevacion@cref}{{[subappendix][1][2147483647,1]A.1}{[1][35][]35}}
\newlabel{sol:eje:1-coterminales}{{A.1}{35}{Tema 1}{section*.11}{}}
\newlabel{sol:eje:1-coterminales@cref}{{[subappendix][1][2147483647,1]A.1}{[1][35][]35}}
\newlabel{sol:eje:1-rope}{{A.1}{35}{Tema 1}{section*.12}{}}
\newlabel{sol:eje:1-rope@cref}{{[subappendix][1][2147483647,1]A.1}{[1][35][]35}}
\newlabel{sol:eje:1-simplificar}{{A.1}{36}{Tema 1}{section*.13}{}}
\newlabel{sol:eje:1-simplificar@cref}{{[subappendix][1][2147483647,1]A.1}{[1][36][]36}}
\newlabel{sol:eje:1-sinx-x-tanx}{{A.1}{36}{Tema 1}{section*.14}{}}
\newlabel{sol:eje:1-sinx-x-tanx@cref}{{[subappendix][1][2147483647,1]A.1}{[1][36][]36}}
\newlabel{sol:eje:1-abs-sin}{{A.1}{36}{Tema 1}{section*.15}{}}
\newlabel{sol:eje:1-abs-sin@cref}{{[subappendix][1][2147483647,1]A.1}{[1][36][]36}}
\newlabel{sol:eje:1-sin-mas-cos}{{A.1}{37}{Tema 1}{section*.16}{}}
\newlabel{sol:eje:1-sin-mas-cos@cref}{{[subappendix][1][2147483647,1]A.1}{[1][37][]37}}
\newlabel{sol:eje:1-periodo-tan}{{A.1}{37}{Tema 1}{section*.17}{}}
\newlabel{sol:eje:1-periodo-tan@cref}{{[subappendix][1][2147483647,1]A.1}{[1][37][]37}}
\newlabel{sol:eje:1-poligono}{{A.1}{37}{Tema 1}{section*.18}{}}
\newlabel{sol:eje:1-poligono@cref}{{[subappendix][1][2147483647,1]A.1}{[1][37][]37}}
\newlabel{sol:eje:2-vf}{{A.2}{38}{Tema 2}{section*.19}{}}
\newlabel{sol:eje:2-vf@cref}{{[subappendix][2][2147483647,1]A.2}{[1][38][]38}}
\newlabel{sol:eje:2-signo-mitad}{{A.2}{38}{Tema 2}{section*.20}{}}
\newlabel{sol:eje:2-signo-mitad@cref}{{[subappendix][2][2147483647,1]A.2}{[1][38][]38}}
\newlabel{sol:eje:2-exactos}{{A.2}{38}{Tema 2}{section*.21}{}}
\newlabel{sol:eje:2-exactos@cref}{{[subappendix][2][2147483647,1]A.2}{[1][38][]38}}
\newlabel{sol:eje:2-cuadrantes}{{A.2}{39}{Tema 2}{section*.22}{}}
\newlabel{sol:eje:2-cuadrantes@cref}{{[subappendix][2][2147483647,1]A.2}{[1][38][]39}}
\newlabel{sol:eje:2-doble}{{A.2}{39}{Tema 2}{section*.23}{}}
\newlabel{sol:eje:2-doble@cref}{{[subappendix][2][2147483647,1]A.2}{[1][39][]39}}
\newlabel{sol:eje:2-pi8}{{A.2}{39}{Tema 2}{section*.24}{}}
\newlabel{sol:eje:2-pi8@cref}{{[subappendix][2][2147483647,1]A.2}{[1][39][]39}}
\newlabel{sol:eje:2-producto}{{A.2}{39}{Tema 2}{section*.25}{}}
\newlabel{sol:eje:2-producto@cref}{{[subappendix][2][2147483647,1]A.2}{[1][39][]39}}
\newlabel{sol:eje:2-potencias}{{A.2}{40}{Tema 2}{section*.26}{}}
\newlabel{sol:eje:2-potencias@cref}{{[subappendix][2][2147483647,1]A.2}{[1][40][]40}}
\newlabel{sol:eje:2-verificar}{{A.2}{40}{Tema 2}{section*.27}{}}
\newlabel{sol:eje:2-verificar@cref}{{[subappendix][2][2147483647,1]A.2}{[1][40][]40}}
\newlabel{sol:eje:2-triple}{{A.2}{40}{Tema 2}{section*.28}{}}
\newlabel{sol:eje:2-triple@cref}{{[subappendix][2][2147483647,1]A.2}{[1][40][]40}}
\newlabel{sol:eje:2-cos36}{{A.2}{41}{Tema 2}{section*.29}{}}
\newlabel{sol:eje:2-cos36@cref}{{[subappendix][2][2147483647,1]A.2}{[1][41][]41}}
\newlabel{sol:eje:2-suma-senos}{{A.2}{41}{Tema 2}{section*.30}{}}
\newlabel{sol:eje:2-suma-senos@cref}{{[subappendix][2][2147483647,1]A.2}{[1][41][]41}}
\newlabel{sol:eje:2-tan-triangulo}{{A.2}{41}{Tema 2}{section*.31}{}}
\newlabel{sol:eje:2-tan-triangulo@cref}{{[subappendix][2][2147483647,1]A.2}{[1][41][]41}}
\newlabel{sol:eje:2-chebyshev}{{A.2}{42}{Tema 2}{section*.32}{}}
\newlabel{sol:eje:2-chebyshev@cref}{{[subappendix][2][2147483647,1]A.2}{[1][42][]42}}
\newlabel{sol:eje:2-potencia-giro}{{A.2}{42}{Tema 2}{section*.33}{}}
\newlabel{sol:eje:2-potencia-giro@cref}{{[subappendix][2][2147483647,1]A.2}{[1][42][]42}}
\newlabel{sol:eje:3-que-teorema}{{A.3}{42}{Tema 3}{section*.34}{}}
\newlabel{sol:eje:3-que-teorema@cref}{{[subappendix][3][2147483647,1]A.3}{[1][42][]42}}
\newlabel{sol:eje:3-imposible}{{A.3}{43}{Tema 3}{section*.35}{}}
\newlabel{sol:eje:3-imposible@cref}{{[subappendix][3][2147483647,1]A.3}{[1][43][]43}}
\newlabel{sol:eje:3-obtuso}{{A.3}{43}{Tema 3}{section*.36}{}}
\newlabel{sol:eje:3-obtuso@cref}{{[subappendix][3][2147483647,1]A.3}{[1][43][]43}}
\newlabel{sol:eje:3-ala}{{A.3}{43}{Tema 3}{section*.37}{}}
\newlabel{sol:eje:3-ala@cref}{{[subappendix][3][2147483647,1]A.3}{[1][43][]43}}
\newlabel{sol:eje:3-ambiguo}{{A.3}{43}{Tema 3}{section*.38}{}}
\newlabel{sol:eje:3-ambiguo@cref}{{[subappendix][3][2147483647,1]A.3}{[1][43][]43}}
\newlabel{sol:eje:3-lal}{{A.3}{44}{Tema 3}{section*.39}{}}
\newlabel{sol:eje:3-lal@cref}{{[subappendix][3][2147483647,1]A.3}{[1][44][]44}}
\newlabel{sol:eje:3-barcos}{{A.3}{44}{Tema 3}{section*.40}{}}
\newlabel{sol:eje:3-barcos@cref}{{[subappendix][3][2147483647,1]A.3}{[1][44][]44}}
\newlabel{sol:eje:3-lll}{{A.3}{44}{Tema 3}{section*.41}{}}
\newlabel{sol:eje:3-lll@cref}{{[subappendix][3][2147483647,1]A.3}{[1][44][]44}}
\newlabel{sol:eje:3-angulo-vectores}{{A.3}{45}{Tema 3}{section*.42}{}}
\newlabel{sol:eje:3-angulo-vectores@cref}{{[subappendix][3][2147483647,1]A.3}{[1][45][]45}}
\newlabel{sol:eje:3-coseno-normas}{{A.3}{45}{Tema 3}{section*.43}{}}
\newlabel{sol:eje:3-coseno-normas@cref}{{[subappendix][3][2147483647,1]A.3}{[1][45][]45}}
\newlabel{sol:eje:3-similitud}{{A.3}{45}{Tema 3}{section*.44}{}}
\newlabel{sol:eje:3-similitud@cref}{{[subappendix][3][2147483647,1]A.3}{[1][45][]45}}
\newlabel{sol:eje:3-circunradio}{{A.3}{46}{Tema 3}{section*.45}{}}
\newlabel{sol:eje:3-circunradio@cref}{{[subappendix][3][2147483647,1]A.3}{[1][46][]46}}
\newlabel{sol:eje:3-heron}{{A.3}{46}{Tema 3}{section*.46}{}}
\newlabel{sol:eje:3-heron@cref}{{[subappendix][3][2147483647,1]A.3}{[1][46][]46}}
\newlabel{sol:eje:3-proyecciones}{{A.3}{46}{Tema 3}{section*.47}{}}
\newlabel{sol:eje:3-proyecciones@cref}{{[subappendix][3][2147483647,1]A.3}{[1][46][]46}}
\newlabel{sol:eje:3-mayor-lado}{{A.3}{47}{Tema 3}{section*.48}{}}
\newlabel{sol:eje:3-mayor-lado@cref}{{[subappendix][3][2147483647,1]A.3}{[1][47][]47}}
\newlabel{sol:eje:4-parametros}{{A.4}{47}{Tema 4}{section*.49}{}}
\newlabel{sol:eje:4-parametros@cref}{{[subappendix][4][2147483647,1]A.4}{[1][47][]47}}
\newlabel{sol:eje:4-dominios}{{A.4}{48}{Tema 4}{section*.50}{}}
\newlabel{sol:eje:4-dominios@cref}{{[subappendix][4][2147483647,1]A.4}{[1][47][]48}}
\newlabel{sol:eje:4-exactos}{{A.4}{48}{Tema 4}{section*.51}{}}
\newlabel{sol:eje:4-exactos@cref}{{[subappendix][4][2147483647,1]A.4}{[1][48][]48}}
\newlabel{sol:eje:4-basicas}{{A.4}{48}{Tema 4}{section*.52}{}}
\newlabel{sol:eje:4-basicas@cref}{{[subappendix][4][2147483647,1]A.4}{[1][48][]48}}
\newlabel{sol:eje:4-cuadratica}{{A.4}{48}{Tema 4}{section*.53}{}}
\newlabel{sol:eje:4-cuadratica@cref}{{[subappendix][4][2147483647,1]A.4}{[1][48][]48}}
\newlabel{sol:eje:4-doble}{{A.4}{49}{Tema 4}{section*.54}{}}
\newlabel{sol:eje:4-doble@cref}{{[subappendix][4][2147483647,1]A.4}{[1][49][]49}}
\newlabel{sol:eje:4-armonico}{{A.4}{49}{Tema 4}{section*.55}{}}
\newlabel{sol:eje:4-armonico@cref}{{[subappendix][4][2147483647,1]A.4}{[1][49][]49}}
\newlabel{sol:eje:4-cuadrado}{{A.4}{49}{Tema 4}{section*.56}{}}
\newlabel{sol:eje:4-cuadrado@cref}{{[subappendix][4][2147483647,1]A.4}{[1][49][]49}}
\newlabel{sol:eje:4-periodo-suma}{{A.4}{50}{Tema 4}{section*.57}{}}
\newlabel{sol:eje:4-periodo-suma@cref}{{[subappendix][4][2147483647,1]A.4}{[1][49][]50}}
\newlabel{sol:eje:4-pe}{{A.4}{50}{Tema 4}{section*.58}{}}
\newlabel{sol:eje:4-pe@cref}{{[subappendix][4][2147483647,1]A.4}{[1][50][]50}}
\newlabel{sol:eje:4-creciente}{{A.4}{50}{Tema 4}{section*.59}{}}
\newlabel{sol:eje:4-creciente@cref}{{[subappendix][4][2147483647,1]A.4}{[1][50][]50}}
\newlabel{sol:eje:4-cos-iguales}{{A.4}{51}{Tema 4}{section*.60}{}}
\newlabel{sol:eje:4-cos-iguales@cref}{{[subappendix][4][2147483647,1]A.4}{[1][51][]51}}
\newlabel{sol:eje:4-arctan-inverso}{{A.4}{51}{Tema 4}{section*.61}{}}
\newlabel{sol:eje:4-arctan-inverso@cref}{{[subappendix][4][2147483647,1]A.4}{[1][51][]51}}
\newlabel{sol:eje:4-machin}{{A.4}{51}{Tema 4}{section*.62}{}}
\newlabel{sol:eje:4-machin@cref}{{[subappendix][4][2147483647,1]A.4}{[1][51][]51}}
\newlabel{sol:eje:4-pe-lineal}{{A.4}{52}{Tema 4}{section*.63}{}}
\newlabel{sol:eje:4-pe-lineal@cref}{{[subappendix][4][2147483647,1]A.4}{[1][51][]52}}
\newlabel{sol:eje:5-partes}{{A.5}{52}{Tema 5}{section*.64}{}}
\newlabel{sol:eje:5-partes@cref}{{[subappendix][5][2147483647,1]A.5}{[1][52][]52}}
\newlabel{sol:eje:5-vf}{{A.5}{52}{Tema 5}{section*.65}{}}
\newlabel{sol:eje:5-vf@cref}{{[subappendix][5][2147483647,1]A.5}{[1][52][]52}}
\newlabel{sol:eje:5-potencias-i}{{A.5}{53}{Tema 5}{section*.66}{}}
\newlabel{sol:eje:5-potencias-i@cref}{{[subappendix][5][2147483647,1]A.5}{[1][52][]53}}
\newlabel{sol:eje:5-operaciones}{{A.5}{53}{Tema 5}{section*.67}{}}
\newlabel{sol:eje:5-operaciones@cref}{{[subappendix][5][2147483647,1]A.5}{[1][53][]53}}
\newlabel{sol:eje:5-inversos}{{A.5}{53}{Tema 5}{section*.68}{}}
\newlabel{sol:eje:5-inversos@cref}{{[subappendix][5][2147483647,1]A.5}{[1][53][]53}}
\newlabel{sol:eje:5-cuadraticas}{{A.5}{53}{Tema 5}{section*.69}{}}
\newlabel{sol:eje:5-cuadraticas@cref}{{[subappendix][5][2147483647,1]A.5}{[1][53][]53}}
\newlabel{sol:eje:5-suma-tripos}{{A.5}{54}{Tema 5}{section*.70}{}}
\newlabel{sol:eje:5-suma-tripos@cref}{{[subappendix][5][2147483647,1]A.5}{[1][54][]54}}
\newlabel{sol:eje:5-raices-cuadradas}{{A.5}{54}{Tema 5}{section*.71}{}}
\newlabel{sol:eje:5-raices-cuadradas@cref}{{[subappendix][5][2147483647,1]A.5}{[1][54][]54}}
\newlabel{sol:eje:5-cuadratica-compleja}{{A.5}{54}{Tema 5}{section*.72}{}}
\newlabel{sol:eje:5-cuadratica-compleja@cref}{{[subappendix][5][2147483647,1]A.5}{[1][54][]54}}
\newlabel{sol:eje:5-lugares}{{A.5}{54}{Tema 5}{section*.73}{}}
\newlabel{sol:eje:5-lugares@cref}{{[subappendix][5][2147483647,1]A.5}{[1][54][]54}}
\newlabel{sol:eje:5-escalar}{{A.5}{55}{Tema 5}{section*.74}{}}
\newlabel{sol:eje:5-escalar@cref}{{[subappendix][5][2147483647,1]A.5}{[1][55][]55}}
\newlabel{sol:eje:5-modulo-producto}{{A.5}{55}{Tema 5}{section*.75}{}}
\newlabel{sol:eje:5-modulo-producto@cref}{{[subappendix][5][2147483647,1]A.5}{[1][55][]55}}
\newlabel{sol:eje:5-igualdad-triangular}{{A.5}{55}{Tema 5}{section*.76}{}}
\newlabel{sol:eje:5-igualdad-triangular@cref}{{[subappendix][5][2147483647,1]A.5}{[1][55][]55}}
\newlabel{sol:eje:5-orden}{{A.5}{56}{Tema 5}{section*.77}{}}
\newlabel{sol:eje:5-orden@cref}{{[subappendix][5][2147483647,1]A.5}{[1][56][]56}}
\newlabel{sol:eje:5-blaschke}{{A.5}{56}{Tema 5}{section*.78}{}}
\newlabel{sol:eje:5-blaschke@cref}{{[subappendix][5][2147483647,1]A.5}{[1][56][]56}}
\newlabel{sol:eje:6-exponenciales}{{A.6}{56}{Tema 6}{section*.79}{}}
\newlabel{sol:eje:6-exponenciales@cref}{{[subappendix][6][2147483647,1]A.6}{[1][56][]56}}
\newlabel{sol:eje:6-vf}{{A.6}{57}{Tema 6}{section*.80}{}}
\newlabel{sol:eje:6-vf@cref}{{[subappendix][6][2147483647,1]A.6}{[1][57][]57}}
\newlabel{sol:eje:6-cuadrante}{{A.6}{57}{Tema 6}{section*.81}{}}
\newlabel{sol:eje:6-cuadrante@cref}{{[subappendix][6][2147483647,1]A.6}{[1][57][]57}}
\newlabel{sol:eje:6-polar}{{A.6}{57}{Tema 6}{section*.82}{}}
\newlabel{sol:eje:6-polar@cref}{{[subappendix][6][2147483647,1]A.6}{[1][57][]57}}
\newlabel{sol:eje:6-producto-polar}{{A.6}{57}{Tema 6}{section*.83}{}}
\newlabel{sol:eje:6-producto-polar@cref}{{[subappendix][6][2147483647,1]A.6}{[1][57][]57}}
\newlabel{sol:eje:6-potencias}{{A.6}{58}{Tema 6}{section*.84}{}}
\newlabel{sol:eje:6-potencias@cref}{{[subappendix][6][2147483647,1]A.6}{[1][58][]58}}
\newlabel{sol:eje:6-tripos}{{A.6}{58}{Tema 6}{section*.85}{}}
\newlabel{sol:eje:6-tripos@cref}{{[subappendix][6][2147483647,1]A.6}{[1][58][]58}}
\newlabel{sol:eje:6-raices}{{A.6}{58}{Tema 6}{section*.86}{}}
\newlabel{sol:eje:6-raices@cref}{{[subappendix][6][2147483647,1]A.6}{[1][58][]58}}
\newlabel{sol:eje:6-polinomios}{{A.6}{59}{Tema 6}{section*.87}{}}
\newlabel{sol:eje:6-polinomios@cref}{{[subappendix][6][2147483647,1]A.6}{[1][58][]59}}
\newlabel{sol:eje:6-giros}{{A.6}{59}{Tema 6}{section*.88}{}}
\newlabel{sol:eje:6-giros@cref}{{[subappendix][6][2147483647,1]A.6}{[1][59][]59}}
\newlabel{sol:eje:6-rope}{{A.6}{59}{Tema 6}{section*.89}{}}
\newlabel{sol:eje:6-rope@cref}{{[subappendix][6][2147483647,1]A.6}{[1][59][]59}}
\newlabel{sol:eje:6-suma-euler}{{A.6}{60}{Tema 6}{section*.90}{}}
\newlabel{sol:eje:6-suma-euler@cref}{{[subappendix][6][2147483647,1]A.6}{[1][59][]60}}
\newlabel{sol:eje:6-cos5}{{A.6}{60}{Tema 6}{section*.91}{}}
\newlabel{sol:eje:6-cos5@cref}{{[subappendix][6][2147483647,1]A.6}{[1][60][]60}}
\newlabel{sol:eje:6-potencias-raices}{{A.6}{60}{Tema 6}{section*.92}{}}
\newlabel{sol:eje:6-potencias-raices@cref}{{[subappendix][6][2147483647,1]A.6}{[1][60][]60}}
\newlabel{sol:eje:6-suma-cosenos}{{A.6}{61}{Tema 6}{section*.93}{}}
\newlabel{sol:eje:6-suma-cosenos@cref}{{[subappendix][6][2147483647,1]A.6}{[1][60][]61}}
\newlabel{sol:eje:6-autovalores}{{A.6}{61}{Tema 6}{section*.94}{}}
\newlabel{sol:eje:6-autovalores@cref}{{[subappendix][6][2147483647,1]A.6}{[1][61][]61}}
\makeatother

\begin{document}
\frontmatter
\begin{titlepage}
\centering
\vspace*{2.5cm}
{\large\scshape Anatomía de un LLM}\par
\vspace{0.4cm}
{\normalsize Fundamentos matemáticos de los modelos de lenguaje}\par
\vspace{2.2cm}
\rule{\textwidth}{0.8pt}\par\vspace{0.5cm}
{\small\scshape Fase 0 \textperiodcentered\ Preliminares matemáticos}\par\vspace{0.3cm}
{\Huge\bfseries Módulo 0.4\par}\vspace{0.3cm}
{\LARGE Trigonometría y números complejos\par}
\vspace{0.5cm}\rule{\textwidth}{0.4pt}\par
\vspace{1.6cm}
{\large Resumen del módulo}\par
\vfill
{\small Edición de octubre de 2026}\par
\end{titlepage}

\chapter*{Sobre este resumen}
\addcontentsline{toc}{chapter}{Sobre este resumen}
\markboth{Sobre este resumen}{}

Este documento reúne lo indispensable de los apuntes del Módulo 0.4: el texto expositivo, todas las definiciones y todos los resultados (teoremas, proposiciones, lemas y corolarios), las tablas y las observaciones de contenido. Se han suprimido las demostraciones, los ejemplos, las lecturas, las figuras, los ejercicios y sus soluciones. La numeración de temas, secciones, enunciados y ecuaciones es la misma que en los apuntes, de modo que cualquier referencia puede buscarse allí; las que remiten a elementos suprimidos (un ejemplo, una figura, una ecuación de una demostración) conservan su número en los apuntes.

\section*{Notación}
\begin{center}
\begin{tabular}{@{}ll@{}}
\toprule
Símbolo & Significado \\
\midrule
$\sin$, $\cos$, $\tan$ & seno, coseno y tangente (en notación escolar, sen, cos, tg) \\
$\arcsin$, $\arccos$, $\arctan$ & funciones trigonométricas inversas (valores principales) \\
$\C$, $i$ & cuerpo de los números complejos; unidad imaginaria, $i^{2}=-1$ \\
$z=a+bi$, $\Real z$, $\Imag z$ & forma binómica; partes real e imaginaria \\
$\bar z$, $\abs{z}$ & conjugado y módulo \\
$\arg z$, $\Arg z$ & un argumento de $z$; argumento principal en $(-\pi,\pi]$ \\
$e^{i\theta}=\cos\theta+i\sin\theta$ & exponencial imaginaria (fórmula de Euler) \\
$R(\alpha)$ & matriz de giro de ángulo $\alpha$ en $\R^{2}$ \\
$u\cdot v$, $\norm{u}$ & producto escalar y norma euclídea \\
\bottomrule
\end{tabular}
\end{center}


\tableofcontents

\mainmatter
% =====================================================================
\chapter{Ángulos, razones trigonométricas y circunferencia unidad}
\label{tema:1}
% =====================================================================

Este tema fija la medida de ángulos en radianes, define las razones trigonométricas primero en el triángulo rectángulo y después, para cualquier número real, mediante la circunferencia unidad, y deduce de las simetrías de la circunferencia las propiedades básicas de signo, paridad y periodicidad. Se sigue el capítulo~7 de~\cite{osat} (§7.1 a §7.4) y el doble enfoque, del triángulo y de la circunferencia unidad, de Stewart, Redlin y Watson~\cite{stewart}.

\section{Ángulos y radianes}

Un ángulo en posición estándar tiene el vértice en el origen y el lado inicial sobre el semieje positivo de abscisas; es positivo si se recorre en sentido antihorario y negativo en sentido horario. Dos ángulos son \emph{coterminales} si tienen el mismo lado final.

\begin{definicion}[Radián]\label{def:radian}
Sea una circunferencia de radio $r>0$ y un ángulo central que abarca un arco de longitud $s$. La medida del ángulo en \emph{radianes} es $\theta=s/r$. Una vuelta completa mide $2\pi$ radianes, de modo que
\[
180^{\circ}=\pi\ \text{rad},\qquad 1^{\circ}=\frac{\pi}{180}\ \text{rad}\approx 0{,}017\,45\ \text{rad},\qquad 1\ \text{rad}=\frac{180^{\circ}}{\pi}\approx 57{,}30^{\circ}.
\]
\end{definicion}

La definición no depende del radio elegido: al multiplicar $r$ por $\lambda$, la semejanza multiplica también $s$ por $\lambda$, y el cociente no cambia. El radián es, por tanto, una razón entre dos longitudes y carece de dimensiones.

\begin{proposicion}[Arco y sector]\label{prop:arco-sector}
En una circunferencia de radio $r$, un ángulo central de $\theta$ radianes ($0\le\theta\le 2\pi$) abarca un arco de longitud $s=r\theta$ y un sector circular de área $A=\tfrac12 r^{2}\theta$.
\end{proposicion}

\addtocounter{definicion}{1}%

\section{Razones trigonométricas en el triángulo rectángulo}

\begin{definicion}[Razones de un ángulo agudo]\label{def:razones-triangulo}
Sea un triángulo rectángulo con un ángulo agudo $\theta$, cateto opuesto $a$, cateto contiguo $b$ e hipotenusa $c$. Se definen
\[
\sin\theta=\frac{a}{c},\qquad \cos\theta=\frac{b}{c},\qquad \tan\theta=\frac{a}{b},
\]
y sus inversas aritméticas $\csc\theta=1/\sin\theta$, $\sec\theta=1/\cos\theta$, $\cot\theta=1/\tan\theta$.
\end{definicion}

\begin{proposicion}\label{prop:razones-bien-definidas}
Las razones de la \cref{def:razones-triangulo} solo dependen de $\theta$, no del triángulo rectángulo elegido.
\end{proposicion}

\begin{proposicion}[Valores notables]\label{prop:notables}
Se cumplen $\sin\frac{\pi}{6}=\cos\frac{\pi}{3}=\frac12$, $\sin\frac{\pi}{3}=\cos\frac{\pi}{6}=\frac{\sqrt3}{2}$ y $\sin\frac{\pi}{4}=\cos\frac{\pi}{4}=\frac{\sqrt2}{2}$.
\end{proposicion}

\section{La circunferencia unidad}

La \cref{def:razones-triangulo} solo tiene sentido para $0<\theta<\pi/2$. Para extenderla a cualquier número real se usa la circunferencia de radio~1 centrada en el origen, $x^{2}+y^{2}=1$ (\cref{fig:circunferencia}).

\begin{definicion}[Seno y coseno de un número real]\label{def:circunferencia}
Para $\theta\in\R$, sea $P(\theta)$ el punto de la circunferencia unidad al que se llega recorriendo sobre ella, desde $(1,0)$, un arco de longitud $\abs{\theta}$, en sentido antihorario si $\theta\ge 0$ y horario si $\theta<0$. Se definen $\cos\theta$ y $\sin\theta$ como las coordenadas de ese punto:
\[
P(\theta)=(\cos\theta,\ \sin\theta).
\]
Cuando $\cos\theta\neq 0$ se define $\tan\theta=\sin\theta/\cos\theta$, y análogamente $\sec$, $\csc$ y $\cot$.
\end{definicion}

\addtocounter{figure}{1}%

\begin{proposicion}\label{prop:coherencia}
Para $0<\theta<\pi/2$, la \cref{def:circunferencia} coincide con la \cref{def:razones-triangulo}.
\end{proposicion}

\begin{teorema}[Propiedades básicas]\label{teo:simetrias}
Para todo $\theta\in\R$ y todo $k\in\Z$:
\begin{enumerate}
  \item (identidad pitagórica) $\sin^{2}\theta+\cos^{2}\theta=1$; en particular, $\abs{\sin\theta}\le 1$ y $\abs{\cos\theta}\le 1$;
  \item (periodicidad) $\sin(\theta+2k\pi)=\sin\theta$ y $\cos(\theta+2k\pi)=\cos\theta$; más aún, $P(\theta)=P(\varphi)$ si y solo si $\theta-\varphi\in2\pi\Z$;
  \item (paridad) $\cos(-\theta)=\cos\theta$ y $\sin(-\theta)=-\sin\theta$;
  \item (suplementario) $\sin(\pi-\theta)=\sin\theta$ y $\cos(\pi-\theta)=-\cos\theta$;
  \item (media vuelta) $\sin(\theta+\pi)=-\sin\theta$ y $\cos(\theta+\pi)=-\cos\theta$;
  \item (complementario) $\sin(\tfrac{\pi}{2}-\theta)=\cos\theta$ y $\cos(\tfrac{\pi}{2}-\theta)=\sin\theta$.
\end{enumerate}
\end{teorema}

\begin{corolario}[Signos por cuadrantes]\label{cor:signos}
El seno es positivo en los cuadrantes I y II; el coseno, en I y IV; la tangente, en I y III.
\end{corolario}

\begin{definicion}[Ángulo de referencia]
El \emph{ángulo de referencia} de $\theta$ es el ángulo agudo $\theta'\in[0,\pi/2]$ que forma el lado final de $\theta$ con el eje $x$. Por el \cref{teo:simetrias}, $\abs{\sin\theta}=\sin\theta'$ y $\abs{\cos\theta}=\cos\theta'$; el signo lo da el \cref{cor:signos}.
\end{definicion}

\addtocounter{definicion}{1}%

\begin{table}[ht]
\centering\small
\caption{Valores notables. Los de los demás cuadrantes se obtienen con el \cref{teo:simetrias}.}\label{tab:notables}
\begin{tabular}{@{}lccccc@{}}
\toprule
$\theta$ & $0$ & $\pi/6$ & $\pi/4$ & $\pi/3$ & $\pi/2$ \\
\midrule
$\sin\theta$ & $0$ & $1/2$ & $\sqrt2/2$ & $\sqrt3/2$ & $1$ \\
$\cos\theta$ & $1$ & $\sqrt3/2$ & $\sqrt2/2$ & $1/2$ & $0$ \\
$\tan\theta$ & $0$ & $\sqrt3/3$ & $1$ & $\sqrt3$ & no definida \\
\bottomrule
\end{tabular}
\end{table}

\addtocounter{definicion}{1}%

\addtocounter{definicion}{1}%


% =====================================================================
\chapter{Identidades trigonométricas}
\label{tema:2}
% =====================================================================

Una identidad trigonométrica es una igualdad entre expresiones trigonométricas que se cumple para todos los valores de la variable en que ambos miembros están definidos. Este tema demuestra las identidades fundamentales a partir de la fórmula del coseno de una diferencia, que a su vez se obtiene de una distancia en la circunferencia unidad, y muestra cómo usarlas para simplificar y verificar expresiones. Se sigue el capítulo~9 de~\cite{osat} (§9.1 a §9.4), equivalente al capítulo~7 de~\cite{osprecalc}; en el \cref{tema:6} se verá una segunda demostración de las fórmulas de la suma, mediante la exponencial compleja, que es la que propone el curso 18.03 del MIT~\cite[§6.4]{mit1803}.

\section{Identidades pitagóricas y verificación de identidades}

\begin{proposicion}[Identidades pitagóricas]\label{prop:pitagoricas}
Para todo $\theta$ en que estén definidas las funciones,
\[
\sin^{2}\theta+\cos^{2}\theta=1,\qquad 1+\tan^{2}\theta=\sec^{2}\theta,\qquad 1+\cot^{2}\theta=\csc^{2}\theta.
\]
\end{proposicion}

Para \emph{verificar} una identidad se transforma uno de sus miembros, mediante identidades ya demostradas y operaciones algebraicas reversibles, hasta obtener el otro. No es válido manipular la igualdad que se quiere probar como si fuera cierta (por ejemplo, elevando ambos miembros al cuadrado), porque así se pueden «demostrar» igualdades falsas~\cite[§9.1]{osat}.

\addtocounter{definicion}{1}%

\section{Fórmulas de la suma y de la diferencia}

\begin{teorema}[Coseno de una diferencia]\label{teo:cos-diferencia}
Para todos $\alpha,\beta\in\R$,
\[
\cos(\alpha-\beta)=\cos\alpha\cos\beta+\sin\alpha\sin\beta .
\]
\end{teorema}

\begin{corolario}[Fórmulas de la suma y la diferencia]\label{cor:suma}
Para todos $\alpha,\beta\in\R$,
\begin{align*}
\cos(\alpha\pm\beta)&=\cos\alpha\cos\beta\mp\sin\alpha\sin\beta,\\
\sin(\alpha\pm\beta)&=\sin\alpha\cos\beta\pm\cos\alpha\sin\beta,\\
\tan(\alpha\pm\beta)&=\frac{\tan\alpha\pm\tan\beta}{1\mp\tan\alpha\tan\beta}\quad\text{(cuando todas las tangentes están definidas).}
\end{align*}
\end{corolario}

\addtocounter{definicion}{1}%

\begin{corolario}[Composición de giros]\label{cor:giros}
Sea $R(\alpha)=\begin{pmatrix}\cos\alpha&-\sin\alpha\\ \sin\alpha&\cos\alpha\end{pmatrix}$. Entonces $R(\alpha)R(\beta)=R(\alpha+\beta)$ para todos $\alpha,\beta$.
\end{corolario}

La matriz $R(\alpha)$ lleva el vector $(\cos\theta,\sin\theta)$ al vector $(\cos(\theta+\alpha),\sin(\theta+\alpha))$, como se ve aplicando el \cref{cor:suma}: es el giro de ángulo $\alpha$. El \cref{cor:giros} dice que girar $\beta$ y después $\alpha$ equivale a girar $\alpha+\beta$.

\section{Ángulo doble y ángulo mitad}

\begin{proposicion}[Ángulo doble]\label{prop:doble}
Para todo $\alpha$,
\[
\sin 2\alpha=2\sin\alpha\cos\alpha,\qquad \cos 2\alpha=\cos^{2}\alpha-\sin^{2}\alpha=2\cos^{2}\alpha-1=1-2\sin^{2}\alpha,
\]
y $\tan 2\alpha=\dfrac{2\tan\alpha}{1-\tan^{2}\alpha}$ cuando los dos miembros están definidos.
\end{proposicion}

\begin{corolario}[Reducción de potencias y ángulo mitad]\label{cor:mitad}
Para todo $\alpha$,
\[
\cos^{2}\alpha=\frac{1+\cos 2\alpha}{2},\qquad \sin^{2}\alpha=\frac{1-\cos 2\alpha}{2}.
\]
En consecuencia, $\cos\frac{\alpha}{2}=\pm\sqrt{\frac{1+\cos\alpha}{2}}$ y $\sin\frac{\alpha}{2}=\pm\sqrt{\frac{1-\cos\alpha}{2}}$, donde el signo es el de la función en el cuadrante de $\alpha/2$, y $\tan\frac{\alpha}{2}=\frac{1-\cos\alpha}{\sin\alpha}=\frac{\sin\alpha}{1+\cos\alpha}$ cuando los denominadores no se anulan.
\end{corolario}

\addtocounter{definicion}{1}%

\section{Fórmulas de producto a suma y de suma a producto}

\begin{proposicion}[Producto a suma]\label{prop:producto-suma}
Para todos $\alpha,\beta$,
\begin{align*}
\cos\alpha\cos\beta&=\tfrac12\bigl[\cos(\alpha-\beta)+\cos(\alpha+\beta)\bigr],\\
\sin\alpha\sin\beta&=\tfrac12\bigl[\cos(\alpha-\beta)-\cos(\alpha+\beta)\bigr],\\
\sin\alpha\cos\beta&=\tfrac12\bigl[\sin(\alpha+\beta)+\sin(\alpha-\beta)\bigr].
\end{align*}
\end{proposicion}

\begin{corolario}[Suma a producto]\label{cor:suma-producto}
Para todos $x,y$,
\begin{align*}
\sin x+\sin y&=2\sin\tfrac{x+y}{2}\cos\tfrac{x-y}{2}, & \sin x-\sin y&=2\cos\tfrac{x+y}{2}\sin\tfrac{x-y}{2},\\
\cos x+\cos y&=2\cos\tfrac{x+y}{2}\cos\tfrac{x-y}{2}, & \cos x-\cos y&=-2\sin\tfrac{x+y}{2}\sin\tfrac{x-y}{2}.
\end{align*}
\end{corolario}

\addtocounter{definicion}{1}%

\addtocounter{definicion}{1}%


% =====================================================================
\chapter{Resolución de triángulos y producto escalar}
\label{tema:3}
% =====================================================================

Este tema extiende la trigonometría a triángulos cualesquiera mediante los teoremas del seno y del coseno, clasifica los casos de resolución de triángulos (incluido el caso ambiguo) y muestra que el teorema del coseno es la misma afirmación que la fórmula geométrica del producto escalar. Se sigue~\cite[§10.1, §10.2 y §10.8]{osat}, el tratamiento de los triángulos oblicuos de~\cite{stewart} y, para el producto escalar, la sección~1.2 de Strang~\cite{strang}.

\section{Área y teorema del seno}

En lo que sigue, $ABC$ es un triángulo con ángulos $A,B,C$ (que se identifican con sus medidas en radianes) y lados $a=BC$, $b=CA$, $c=AB$, de modo que cada lado es el opuesto al ángulo de la misma letra. Se cumple $A+B+C=\pi$ y todos los ángulos están en $(0,\pi)$, de modo que sus senos son positivos.

\begin{proposicion}[Área]\label{prop:area}
El área de $ABC$ es $S=\tfrac12 ab\sin C=\tfrac12 bc\sin A=\tfrac12 ca\sin B$.
\end{proposicion}

\begin{teorema}[Teorema del seno]\label{teo:seno}
En todo triángulo,
\[
\frac{a}{\sin A}=\frac{b}{\sin B}=\frac{c}{\sin C}.
\]
\end{teorema}

El valor común de los tres cocientes es el diámetro $2R$ de la circunferencia circunscrita (\cref{eje:3-circunradio}).

\section{Teorema del coseno}

\begin{teorema}[Teorema del coseno]\label{teo:coseno}
En todo triángulo, $c^{2}=a^{2}+b^{2}-2ab\cos C$, y análogamente para los otros dos lados.
\end{teorema}

\begin{corolario}\label{cor:pitagoras}
$C=\pi/2$ si y solo si $c^{2}=a^{2}+b^{2}$ (teorema de Pitágoras y su recíproco); $C$ es obtuso si y solo si $c^{2}>a^{2}+b^{2}$, y agudo si y solo si $c^{2}<a^{2}+b^{2}$.
\end{corolario}

\section{Resolución de triángulos}

\emph{Resolver} un triángulo es calcular sus seis elementos a partir de tres datos que incluyan al menos un lado. La \cref{tab:casos} resume qué teorema se usa en cada caso~\cite[§10.1--10.2]{osat}.

\begin{table}[ht]
\centering\small
\caption{Casos de resolución de triángulos.}\label{tab:casos}
\begin{tabular}{@{}>{\raggedright\arraybackslash}p{4.6cm}>{\raggedright\arraybackslash}p{4.6cm}>{\raggedright\arraybackslash}p{4.6cm}@{}}
\toprule
Datos & Primer paso & Número de soluciones \\
\midrule
LLL (tres lados) & teorema del coseno para un ángulo & una, si se cumple la desigualdad triangular \\
LAL (dos lados y el ángulo comprendido) & teorema del coseno para el tercer lado & una \\
ALA o AAL (dos ángulos y un lado) & $C=\pi-A-B$ y teorema del seno & una, si $A+B<\pi$ \\
LLA (dos lados y un ángulo no comprendido) & teorema del seno & cero, una o dos (caso ambiguo) \\
\bottomrule
\end{tabular}
\end{table}

\begin{proposicion}[Caso ambiguo]\label{prop:ambiguo}
Se dan $a$, $b$ y el ángulo $A$ opuesto a $a$, con $A$ agudo, y sea $h=b\sin A$. Entonces existen: ningún triángulo si $a<h$; uno (rectángulo en $B$) si $a=h$; dos si $h<a<b$; y uno si $a\ge b$.
\end{proposicion}

\addtocounter{definicion}{1}%

Cuando se conocen los tres lados, el área puede calcularse sin ángulos mediante la \emph{fórmula de Herón}, $S=\sqrt{s(s-a)(s-b)(s-c)}$ con $s=\frac{a+b+c}{2}$, que se demuestra en el \cref{eje:3-heron}.

\section{Producto escalar y ángulo entre vectores}

\begin{definicion}[Producto escalar]
Para $u=(u_1,\dots,u_n)$ y $v=(v_1,\dots,v_n)$ en $\R^{n}$, el \emph{producto escalar} es $u\cdot v=\sum_{i=1}^{n}u_iv_i$ y la \emph{norma} es $\norm{u}=\sqrt{u\cdot u}$.
\end{definicion}

El producto escalar es simétrico y bilineal; en particular,
\begin{equation}\label{eq:norma-diferencia}
\norm{u-v}^{2}=(u-v)\cdot(u-v)=\norm{u}^{2}+\norm{v}^{2}-2\,u\cdot v .
\end{equation}

\begin{teorema}[Forma geométrica del producto escalar]\label{teo:escalar}
Sean $u,v\in\R^{2}$ (o $\R^{3}$) no nulos y $\varphi\in[0,\pi]$ el ángulo que forman. Entonces
\[
u\cdot v=\norm{u}\,\norm{v}\cos\varphi .
\]
\end{teorema}

\begin{corolario}\label{cor:escalar}
Para $u,v$ no nulos del plano o del espacio: (i) $u\perp v$ si y solo si $u\cdot v=0$; (ii) $\abs{u\cdot v}\le\norm{u}\norm{v}$ (desigualdad de Cauchy--Schwarz); (iii) la componente de $u$ en la dirección de $v$ es $\norm{u}\cos\varphi=u\cdot v/\norm{v}$, y la proyección de $u$ sobre la recta de $v$ es $\frac{u\cdot v}{\norm{v}^{2}}\,v$.
\end{corolario}

En $\R^{n}$ con $n>3$ no hay triángulos que dibujar y se procede al revés: la desigualdad de Cauchy--Schwarz se demuestra algebraicamente y se \emph{define} el ángulo por $\cos\varphi=u\cdot v/(\norm{u}\norm{v})$ (módulo~2, tema~2). El cociente recibe el nombre de \emph{similitud coseno}.

\addtocounter{definicion}{1}%

\addtocounter{definicion}{1}%

\addtocounter{definicion}{1}%


% =====================================================================
\chapter{Funciones trigonométricas, inversas y ecuaciones}
\label{tema:4}
% =====================================================================

Este tema estudia el seno, el coseno y la tangente como funciones de una variable real: su periodicidad, sus gráficas y las transformaciones que producen las sinusoides generales. Después restringe sus dominios para obtener funciones inversas y, con ellas, describe todas las soluciones de las ecuaciones trigonométricas. Se sigue el capítulo~8 de~\cite{osat} (§8.1 a §8.3) y la §9.5 del mismo libro; la conexión con la codificación de la posición procede de Vaswani \emph{et al.}~\cite[§3.5]{vaswani}.

\section{Funciones periódicas}

\begin{definicion}[Función periódica]
Una función $f\colon D\to\R$ es \emph{periódica} de periodo $T>0$ si, para todo $x\in D$, se cumple $x+T\in D$, $x-T\in D$ y $f(x+T)=f(x)$. Si existe un periodo mínimo, se llama \emph{periodo fundamental}.
\end{definicion}

\begin{proposicion}\label{prop:periodos}
(i) El seno y el coseno tienen periodo fundamental $2\pi$; la tangente, $\pi$. (ii) Si $f$ tiene periodo fundamental $T$ y $\omega>0$, la función $g(x)=f(\omega x)$ tiene periodo fundamental $T/\omega$.
\end{proposicion}

\section{Gráficas y sinusoides}

La \cref{fig:graficas} muestra las gráficas del seno y del coseno. Ambas oscilan entre $-1$ y $1$, y la del coseno es la del seno desplazada $\pi/2$ hacia la izquierda, porque $\cos x=\sin(x+\frac{\pi}{2})$ (\cref{teo:simetrias}(vi) con la paridad). La tangente tiene periodo $\pi$, ceros en $k\pi$ y asíntotas verticales en $\frac{\pi}{2}+k\pi$, donde se anula el coseno.

\addtocounter{figure}{1}%

\begin{definicion}[Sinusoide]
Una \emph{sinusoide} es una función $y=A\sin\bigl(\omega(x-x_0)\bigr)+D$ con $A\ne0$ y $\omega>0$. Se llaman \emph{amplitud} a $\abs{A}$, \emph{periodo} a $2\pi/\omega$, \emph{frecuencia angular} a $\omega$, \emph{desfase} (desplazamiento horizontal) a $x_0$ y \emph{desplazamiento vertical} a $D$.
\end{definicion}

Por la \cref{prop:periodos}(ii), el periodo de la sinusoide es efectivamente $2\pi/\omega$. Al escribir $y=A\sin(\omega x-c)+D$, el desfase es $x_0=c/\omega$, no $c$.

\begin{proposicion}[Suma de un seno y un coseno]\label{prop:armonico}
Sean $a,b\in\R$ no ambos nulos y $R=\sqrt{a^{2}+b^{2}}$. Existe $\varphi\in(-\pi,\pi]$ con $\cos\varphi=a/R$ y $\sin\varphi=b/R$, y para todo $x$
\[
a\sin x+b\cos x=R\sin(x+\varphi).
\]
En particular, el máximo de $a\sin x+b\cos x$ es $R$ y el mínimo, $-R$.
\end{proposicion}

\section{Funciones trigonométricas inversas}

Una función periódica no es inyectiva, luego no tiene inversa. Se restringe a un intervalo donde sea biyectiva sobre su imagen~\cite[§8.3]{osat}.

\begin{definicion}[Inversas principales]\label{def:inversas}
\begin{enumerate}
  \item $\sin\colon[-\frac{\pi}{2},\frac{\pi}{2}]\to[-1,1]$ es biyectiva; su inversa es $\arcsin\colon[-1,1]\to[-\frac{\pi}{2},\frac{\pi}{2}]$.
  \item $\cos\colon[0,\pi]\to[-1,1]$ es biyectiva; su inversa es $\arccos\colon[-1,1]\to[0,\pi]$.
  \item $\tan\colon(-\frac{\pi}{2},\frac{\pi}{2})\to\R$ es biyectiva; su inversa es $\arctan\colon\R\to(-\frac{\pi}{2},\frac{\pi}{2})$.
\end{enumerate}
\end{definicion}

La inyectividad de estas restricciones se demuestra en el \cref{eje:4-creciente} (el seno es estrictamente creciente en $[-\frac{\pi}{2},\frac{\pi}{2}]$) y de forma análoga para las otras dos. La sobreyectividad es geométrica: cada recta horizontal $y=t$ con $\abs{t}\le1$ corta la semicircunferencia derecha en un punto, y cada recta vertical $x=t$ corta la superior en un punto; para la tangente, cada recta por el origen no vertical corta la semicircunferencia derecha abierta.

\begin{proposicion}[Propiedades de las inversas]\label{prop:inversas}
\begin{enumerate}
  \item $\sin(\arcsin x)=x$ para $x\in[-1,1]$, y $\arcsin(\sin\theta)=\theta$ si y solo si $\theta\in[-\frac{\pi}{2},\frac{\pi}{2}]$.
  \item $\cos(\arcsin x)=\sqrt{1-x^{2}}$ y $\sin(\arccos x)=\sqrt{1-x^{2}}$ para $x\in[-1,1]$.
  \item $\arcsin x+\arccos x=\frac{\pi}{2}$ para $x\in[-1,1]$.
  \item $\arcsin$ y $\arctan$ son impares, y $\arccos(-x)=\pi-\arccos x$.
\end{enumerate}
\end{proposicion}

\addtocounter{definicion}{1}%

\section{Ecuaciones trigonométricas}

\begin{teorema}[Soluciones generales]\label{teo:ecuaciones}
Sea $a\in\R$.
\begin{enumerate}
  \item Si $\abs{a}\le1$ y $\alpha=\arcsin a$, entonces $\sin x=a$ si y solo si $x=\alpha+2k\pi$ o $x=\pi-\alpha+2k\pi$ con $k\in\Z$.
  \item Si $\abs{a}\le1$ y $\beta=\arccos a$, entonces $\cos x=a$ si y solo si $x=\pm\beta+2k\pi$ con $k\in\Z$.
  \item Para todo $a$, $\tan x=a$ si y solo si $x=\arctan a+k\pi$ con $k\in\Z$.
\end{enumerate}
Si $\abs{a}>1$, las ecuaciones $\sin x=a$ y $\cos x=a$ no tienen solución.
\end{teorema}

\addtocounter{definicion}{1}%

\addtocounter{definicion}{1}%

\addtocounter{definicion}{1}%


% =====================================================================
\chapter{Números complejos en forma binómica}
\label{tema:5}
% =====================================================================

Los números complejos amplían los reales con una raíz cuadrada de $-1$. Con ellos toda ecuación polinómica tiene solución, y las operaciones de giro y homotecia del plano se escriben como productos. Este tema construye $\C$, demuestra las propiedades del conjugado y del módulo y resuelve ecuaciones de segundo grado con coeficientes complejos. Se sigue la §2.4 de~\cite{osat}, la sección~5 de las notas de 18.03 del MIT~\cite{mit1803} y la §1.1 de las notas de «Vectors and Matrices» de Cambridge~\cite{cambvm}.

\section{El cuerpo de los números complejos}

\begin{definicion}[Números complejos]\label{def:complejos}
Un \emph{número complejo} es una expresión $z=a+bi$ con $a,b\in\R$; se llaman \emph{parte real} $\Real z=a$ y \emph{parte imaginaria} $\Imag z=b$. Dos complejos son iguales si tienen iguales las partes reales y las imaginarias. Con $z=a+bi$ y $w=c+di$ se definen
\[
z+w=(a+c)+(b+d)i,\qquad zw=(ac-bd)+(ad+bc)i .
\]
El conjunto de los números complejos se denota $\C$. Se identifica el real $a$ con $a+0i$, y se escribe $i=0+1i$.
\end{definicion}

La regla del producto es la que resulta de operar como con polinomios en $i$ y sustituir $i^{2}=-1$: en efecto, con la definición, $i\cdot i=(0\cdot0-1\cdot1)+(0\cdot1+1\cdot0)i=-1$. Geométricamente, $a+bi$ se representa como el punto $(a,b)$ del \emph{plano complejo}, con el eje real horizontal y el imaginario vertical (\cref{fig:argand}).

\begin{teorema}\label{teo:cuerpo}
Con estas operaciones, $\C$ es un cuerpo: la suma y el producto son asociativos y conmutativos, el producto es distributivo respecto de la suma, $0$ y $1$ son los elementos neutros, y todo $z=a+bi\ne0$ tiene inverso
\[
z^{-1}=\frac{a}{a^{2}+b^{2}}-\frac{b}{a^{2}+b^{2}}\,i .
\]
\end{teorema}

\addtocounter{definicion}{1}%

\section{El conjugado}

\begin{definicion}[Conjugado]
El \emph{conjugado} de $z=a+bi$ es $\bar z=a-bi$. Geométricamente, es la reflexión de $z$ en el eje real.
\end{definicion}

\begin{proposicion}[Propiedades del conjugado]\label{prop:conjugado}
Para todos $z,w\in\C$:
\begin{enumerate}
  \item $\overline{z+w}=\bar z+\bar w$, \quad $\overline{zw}=\bar z\,\bar w$, \quad $\overline{z/w}=\bar z/\bar w$ si $w\ne0$, \quad $\bar{\bar z}=z$;
  \item $z+\bar z=2\Real z$ y $z-\bar z=2i\Imag z$;
  \item $z$ es real si y solo si $z=\bar z$;
  \item $z\bar z=(\Real z)^{2}+(\Imag z)^{2}$, que es un número real mayor o igual que~0.
\end{enumerate}
\end{proposicion}

La propiedad~(iv) da el método práctico para dividir: se multiplican numerador y denominador por el conjugado del denominador,
\[
\frac{z}{w}=\frac{z\,\bar w}{w\,\bar w}=\frac{z\,\bar w}{c^{2}+d^{2}} .
\]

\addtocounter{definicion}{1}%

\section{El módulo}

\begin{definicion}[Módulo]
El \emph{módulo} de $z=a+bi$ es $\abs{z}=\sqrt{a^{2}+b^{2}}=\sqrt{z\bar z}$, la distancia de $z$ al origen. La distancia entre $z$ y $w$ es $\abs{z-w}$.
\end{definicion}

\begin{proposicion}[Propiedades del módulo]\label{prop:modulo}
Para todos $z,w\in\C$:
\begin{enumerate}
  \item $\abs{z}\ge0$, con igualdad si y solo si $z=0$; \quad $\abs{\bar z}=\abs{z}$;
  \item $\abs{\Real z}\le\abs{z}$ y $\abs{\Imag z}\le\abs{z}$;
  \item $\abs{zw}=\abs{z}\abs{w}$ y, si $w\ne0$, $\abs{z/w}=\abs{z}/\abs{w}$;
  \item (desigualdad triangular) $\abs{z+w}\le\abs{z}+\abs{w}$;
  \item $\bigl|\abs{z}-\abs{w}\bigr|\le\abs{z-w}$.
\end{enumerate}
\end{proposicion}

\section{Ecuaciones polinómicas}

\begin{proposicion}[Raíces cuadradas]\label{prop:raiz-cuadrada}
Sea $w=a+bi$ con $b\ne0$. Las soluciones de $z^{2}=w$ son $z=\pm(x+yi)$ con
\[
x=\sqrt{\frac{\abs{w}+a}{2}},\qquad y=\operatorname{signo}(b)\sqrt{\frac{\abs{w}-a}{2}} .
\]
\end{proposicion}

\addtocounter{definicion}{1}%

\begin{proposicion}[Fórmula cuadrática en $\C$]\label{prop:cuadratica}
Sean $a,b,c\in\C$ con $a\ne0$, y sea $\delta\in\C$ con $\delta^{2}=b^{2}-4ac$. Las soluciones de $az^{2}+bz+c=0$ son $z=\dfrac{-b\pm\delta}{2a}$.
\end{proposicion}

\begin{proposicion}[Raíces conjugadas]\label{prop:conjugadas}
Si $p(z)=\sum_{k=0}^{n}c_kz^{k}$ tiene coeficientes reales, entonces $p(\bar z)=\overline{p(z)}$ para todo $z$. En particular, si $z$ es raíz de $p$, también lo es $\bar z$.
\end{proposicion}

\begin{teorema}[Teorema fundamental del álgebra]\label{teo:fundamental}
Todo polinomio de grado $n\ge1$ con coeficientes complejos tiene al menos una raíz en $\C$. En consecuencia, se factoriza como $c\prod_{k=1}^{n}(z-z_k)$ y tiene exactamente $n$ raíces contadas con multiplicidad.
\end{teorema}

La demostración requiere herramientas de análisis y se omite; véanse~\cite[§5.3]{mit1803} y~\cite{needham}. La consecuencia se obtiene por inducción sobre $n$, dividiendo por $z-z_1$ cada vez que se encuentra una raíz.

\section{El plano complejo y el producto escalar}

\addtocounter{figure}{1}%

Al identificar $\C$ con $\R^{2}$, la suma de complejos es la suma de vectores, y el producto por un real $\lambda$ es el producto de un vector por un escalar. El producto de dos complejos no tiene análogo en $\R^{2}$ y es el objeto del \cref{tema:6}. El conjugado permite recuperar el producto escalar y el área orientada:

\begin{proposicion}\label{prop:escalar-complejo}
Sean $z=a+bi$ y $w=c+di$, identificados con $u=(a,b)$ y $v=(c,d)$. Entonces
\[
\Real(z\bar w)=ac+bd=u\cdot v,\qquad \Imag(\bar z w)=ad-bc=\det\begin{pmatrix}a&c\\ b&d\end{pmatrix}.
\]
\end{proposicion}

\addtocounter{definicion}{1}%

\addtocounter{definicion}{1}%


% =====================================================================
\chapter{Forma polar, fórmula de Euler, De Moivre y rotaciones}
\label{tema:6}
% =====================================================================

La forma binómica es cómoda para sumar; la forma polar, para multiplicar. Este tema demuestra que al multiplicar dos complejos se multiplican los módulos y se suman los argumentos, introduce la exponencial imaginaria $e^{i\theta}$ y la fórmula de Euler, deduce el teorema de De Moivre y la fórmula de las raíces $n$-ésimas, y termina interpretando la multiplicación por $e^{i\alpha}$ como el giro de ángulo $\alpha$ del plano, que es la idea sobre la que se construye RoPE. Se siguen la sección~6 de las notas de 18.03 del MIT~\cite{mit1803}, las §1.2 a §1.5 de las notas de Cambridge~\cite{cambvm}, la §10.5 de~\cite{osat} y la lectura geométrica del capítulo~1 de Needham~\cite{needham}.

\section{Argumento y forma polar}

\begin{definicion}[Argumento y forma polar]\label{def:polar}
Sea $z\ne0$ y $r=\abs{z}$. Un \emph{argumento} de $z$ es cualquier $\theta\in\R$ tal que
\[
z=r(\cos\theta+i\sin\theta).
\]
Esta expresión es la \emph{forma polar} de $z$. El \emph{argumento principal} $\Arg z$ es el único argumento en $(-\pi,\pi]$. El número $0$ no tiene argumento.
\end{definicion}

Todo $z\ne0$ tiene argumento: $z/r$ tiene módulo~1, luego es un punto $P(\theta)$ de la circunferencia unidad. Por el \cref{teo:simetrias}(ii), dos argumentos de $z$ difieren en un múltiplo entero de $2\pi$, y exactamente uno está en $(-\pi,\pi]$.

\begin{proposicion}[Cálculo del argumento principal]\label{prop:arg}
Sea $z=a+bi\ne0$. Si $a>0$, $\Arg z=\arctan(b/a)$. Si $a<0$, $\Arg z=\arctan(b/a)+\pi$ cuando $b\ge0$ y $\arctan(b/a)-\pi$ cuando $b<0$. Si $a=0$, $\Arg z=\pi/2$ cuando $b>0$ y $-\pi/2$ cuando $b<0$.
\end{proposicion}

\addtocounter{definicion}{1}%

\section{El producto en forma polar}

\begin{teorema}[Producto y cociente en forma polar]\label{teo:producto-polar}
Sean $z_1=r_1(\cos\theta_1+i\sin\theta_1)$ y $z_2=r_2(\cos\theta_2+i\sin\theta_2)$. Entonces
\[
z_1z_2=r_1r_2\bigl(\cos(\theta_1+\theta_2)+i\sin(\theta_1+\theta_2)\bigr),
\]
y, si $z_2\ne0$, $\dfrac{z_1}{z_2}=\dfrac{r_1}{r_2}\bigl(\cos(\theta_1-\theta_2)+i\sin(\theta_1-\theta_2)\bigr)$.
\end{teorema}

\begin{corolario}\label{cor:arg-producto}
Para $z,w\ne0$: $\abs{zw}=\abs{z}\abs{w}$ y $\arg(zw)=\arg z+\arg w$, en el sentido de que la suma de un argumento de $z$ y uno de $w$ es un argumento de $zw$. Para los argumentos principales, la igualdad solo vale salvo un múltiplo de $2\pi$.
\end{corolario}

\emph{Interpretación geométrica.} Fijado $w=\rho(\cos\alpha+i\sin\alpha)$, la aplicación $z\mapsto wz$ multiplica el módulo de cada punto por $\rho$ y suma $\alpha$ a su argumento: es la composición de un giro de ángulo $\alpha$ alrededor del origen y una homotecia de razón $\rho$. Esta lectura, que Needham sitúa en el centro de su exposición~\cite[cap.~1]{needham}, convierte en evidentes muchas identidades: por ejemplo, multiplicar por $i$ es girar un ángulo recto, y por eso $i^{2}=-1$ (dos giros rectos forman media vuelta).

\section{La exponencial imaginaria y la fórmula de Euler}

\begin{definicion}[Exponencial imaginaria]\label{def:euler}
Para $\theta\in\R$ se define
\[
e^{i\theta}=\cos\theta+i\sin\theta,
\]
y, para $z=x+iy\in\C$, $e^{z}=e^{x}e^{iy}=e^{x}(\cos y+i\sin y)$.
\end{definicion}

La notación no es arbitraria. Si se sustituye formalmente $x=i\theta$ en la serie $e^{x}=\sum_{n\ge0}x^{n}/n!$ y se agrupan los términos pares e impares usando $i^{2}=-1$, se obtienen exactamente las series de $\cos\theta$ y de $i\sin\theta$; las notas de Cambridge toman la serie como definición de $e^{z}$ y deducen la fórmula anterior~\cite[§1.2]{cambvm}. Las notas de 18.03 la justifican de otro modo: $f(\theta)=\cos\theta+i\sin\theta$ es la única función con $f'=if$ y $f(0)=1$, la misma ecuación que caracteriza a $e^{i\theta}$~\cite[§6.2]{mit1803}. Ambos argumentos requieren cálculo; aquí basta la definición y sus propiedades algebraicas.

\begin{proposicion}[Propiedades de $e^{i\theta}$]\label{prop:exp}
Para todos $\alpha,\beta,\theta\in\R$ y $n\in\Z$:
\begin{enumerate}
  \item $e^{i\alpha}e^{i\beta}=e^{i(\alpha+\beta)}$ y, más en general, $e^{z}e^{w}=e^{z+w}$ para $z,w\in\C$;
  \item $\abs{e^{i\theta}}=1$ y $\overline{e^{i\theta}}=e^{-i\theta}=1/e^{i\theta}$;
  \item $e^{i\theta}=1$ si y solo si $\theta\in2\pi\Z$; en particular, $e^{i\pi}=-1$ (identidad de Euler) y $e^{2\pi i}=1$;
  \item $\cos\theta=\dfrac{e^{i\theta}+e^{-i\theta}}{2}$ y $\sin\theta=\dfrac{e^{i\theta}-e^{-i\theta}}{2i}$.
\end{enumerate}
\end{proposicion}

Con esta notación, todo $z\ne0$ se escribe en \emph{forma exponencial} $z=re^{i\theta}$, con $r=\abs{z}$ y $\theta$ un argumento, y el \cref{teo:producto-polar} se resume en $r_1e^{i\theta_1}\cdot r_2e^{i\theta_2}=r_1r_2e^{i(\theta_1+\theta_2)}$. El conjugado es $\overline{re^{i\theta}}=re^{-i\theta}$.

\section{El teorema de De Moivre}

\begin{teorema}[De Moivre]\label{teo:demoivre}
Para todo $\theta\in\R$ y todo $n\in\Z$,
\[
(\cos\theta+i\sin\theta)^{n}=\cos n\theta+i\sin n\theta,\qquad\text{es decir,}\qquad (e^{i\theta})^{n}=e^{in\theta}.
\]
\end{teorema}

\begin{corolario}
Si $z=re^{i\theta}$, entonces $z^{n}=r^{n}e^{in\theta}$ para todo $n\in\Z$ (para $n<0$, con $z\ne0$).
\end{corolario}

\addtocounter{definicion}{1}%

\section{Raíces \texorpdfstring{$n$}{n}-ésimas y raíces de la unidad}

\begin{teorema}[Raíces $n$-ésimas]\label{teo:raices}
Sean $w=\rho e^{i\varphi}\ne0$ y $n\ge1$. La ecuación $z^{n}=w$ tiene exactamente $n$ soluciones distintas,
\[
z_k=\rho^{1/n}\,e^{i(\varphi+2\pi k)/n},\qquad k=0,1,\dots,n-1,
\]
donde $\rho^{1/n}$ es la raíz $n$-ésima real positiva de $\rho$.
\end{teorema}

Las $n$ raíces tienen el mismo módulo y argumentos separados $2\pi/n$: son los vértices de un polígono regular de $n$ lados centrado en el origen.

\begin{definicion}[Raíces de la unidad]
Las soluciones de $z^{n}=1$ son las \emph{raíces $n$-ésimas de la unidad}, $1,\omega,\omega^{2},\dots,\omega^{n-1}$ con $\omega=e^{2\pi i/n}$.
\end{definicion}

\begin{proposicion}\label{prop:suma-raices}
Si $n\ge2$, $\;1+\omega+\omega^{2}+\dots+\omega^{n-1}=0$.
\end{proposicion}

Geométricamente, la suma de los vectores que van del centro a los vértices de un polígono regular es nula: el conjunto es invariante por el giro de ángulo $2\pi/n$ (multiplicar la suma por $\omega$ la deja igual), y el único vector fijo por un giro no trivial es el cero.

\addtocounter{figure}{1}%

\section{Giros del plano como multiplicación por \texorpdfstring{$e^{i\alpha}$}{e\textasciicircum ia}}

\begin{proposicion}[Giros]\label{prop:giro-complejo}
Identifíquese $z=x+iy$ con el vector $(x,y)\T$. Para todo $\alpha\in\R$:
\begin{enumerate}
  \item la aplicación $z\mapsto e^{i\alpha}z$ corresponde a $(x,y)\T\mapsto R(\alpha)(x,y)\T$, donde $R(\alpha)$ es la matriz del \cref{cor:giros};
  \item conserva los módulos y los productos escalares: $\abs{e^{i\alpha}z}=\abs{z}$ y $\Real\bigl(e^{i\alpha}z\,\overline{e^{i\alpha}w}\bigr)=\Real(z\bar w)$;
  \item es biyectiva, con inversa $z\mapsto e^{-i\alpha}z$, y la composición de los giros de ángulos $\alpha$ y $\beta$ es el giro de ángulo $\alpha+\beta$.
\end{enumerate}
\end{proposicion}

La \cref{prop:giro-complejo}(iii) es la versión compleja del \cref{cor:giros}, y su demostración es más corta porque la regla de los exponentes ya contiene las fórmulas de la suma. Para girar alrededor de un punto $c$ distinto del origen se traslada, se gira y se deshace la traslación: $z\mapsto c+e^{i\alpha}(z-c)$.

\begin{corolario}[Productos escalares de vectores girados]\label{cor:rope}
Sean $z_1,z_2\in\C$, $\theta\in\R$ y $m,n\in\Z$. Entonces
\[
\Real\bigl(z_1e^{im\theta}\,\overline{z_2e^{in\theta}}\bigr)=\Real\bigl(z_1\bar z_2\,e^{i(m-n)\theta}\bigr),
\]
que solo depende de $m$ y $n$ a través de la diferencia $m-n$.
\end{corolario}

\addtocounter{definicion}{1}%

\addtocounter{definicion}{1}%



\backmatter
% =====================================================================
\begin{thebibliography}{99}
\addcontentsline{toc}{chapter}{Bibliografía}
\markboth{Bibliografía}{}
\small

\bibitem{osat}
J.~Abramson \emph{et al.}
\newblock \emph{Algebra and Trigonometry}, 2.ª ed.
\newblock OpenStax, Rice University, 2021.
\newblock §2.4 «Complex Numbers»; cap.~7 «The Unit Circle: Sine and Cosine Functions» (§7.1--7.4); cap.~8 «Periodic Functions» (§8.1--8.3); cap.~9 «Trigonometric Identities and Equations» (§9.1--9.5); cap.~10 «Further Applications of Trigonometry» (§10.1, §10.2, §10.5, §10.8).
\newblock \url{https://openstax.org/details/books/algebra-and-trigonometry-2e}

\bibitem{osprecalc}
J.~Abramson \emph{et al.}
\newblock \emph{Precalculus}, 2.ª ed.
\newblock OpenStax, Rice University, 2021.
\newblock Cap.~7 «Trigonometric Identities and Equations» (§7.1--7.6) y §8.5 «Polar Form of Complex Numbers».
\newblock \url{https://openstax.org/details/books/precalculus-2e}

\bibitem{stewart}
J.~Stewart, L.~Redlin y S.~Watson.
\newblock \emph{Precalculus: Mathematics for Calculus}, 7.ª ed.
\newblock Cengage Learning, 2016.
\newblock Capítulos de trigonometría (enfoque de la circunferencia unidad y del triángulo rectángulo, trigonometría analítica, coordenadas polares y números complejos).

\bibitem{mit1803}
H.~R. Miller.
\newblock \emph{18.03 Supplementary Notes}.
\newblock Massachusetts Institute of Technology, MIT OpenCourseWare, 18.03 Differential Equations, primavera de 2010.
\newblock §5 «The algebra of complex numbers» (5.1--5.3) y §6 «The complex exponential» (6.1--6.5).
\newblock \url{https://ocw.mit.edu/courses/18-03-differential-equations-spring-2010/pages/readings/notes_exe/}

\bibitem{cambvm}
N.~Peake (profesor); D.~Chua (transcripción).
\newblock \emph{Part IA Vectors and Matrices}, apuntes de clase.
\newblock University of Cambridge, Mathematical Tripos, Michaelmas 2014.
\newblock Cap.~1 «Complex numbers» (§1.1--1.6).
\newblock \url{https://dec41.user.srcf.net/h/IA_M/vectors_and_matrices}

\bibitem{tripos2013}
University of Cambridge.
\newblock \emph{Mathematical Tripos Part IA}, Paper~1, 30 de mayo de 2013.
\newblock Pregunta~1C (Vectors and Matrices).
\newblock \url{https://www.maths.cam.ac.uk/undergrad/pastpapers/files/2013/PaperIA_1.pdf}

\bibitem{tripos2019}
University of Cambridge.
\newblock \emph{Mathematical Tripos Part IA}, Paper~1, 2019.
\newblock Pregunta~1C (Vectors and Matrices).
\newblock \url{https://www.maths.cam.ac.uk/undergrad/pastpapers/files/2019/paperia_1_2019.pdf}

\bibitem{needham}
T.~Needham.
\newblock \emph{Visual Complex Analysis}, edición del 25.º aniversario.
\newblock Oxford University Press, 2023 (1.ª ed., 1997).
\newblock Cap.~1 «Geometry and Complex Arithmetic».

\bibitem{strang}
G.~Strang.
\newblock \emph{Introduction to Linear Algebra}, 4.ª ed.
\newblock Wellesley-Cambridge Press, 2009.
\newblock Texto de los cursos 18.06 y 18.06SC del MIT; §1.2 «Lengths and Dot Products».

\bibitem{roformer}
J.~Su, Y.~Lu, S.~Pan, A.~Murtadha, B.~Wen e Y.~Liu.
\newblock RoFormer: Enhanced Transformer with rotary position embedding.
\newblock Prepublicación, 2021. Publicado en \emph{Neurocomputing}, 568, 2024.
\newblock \url{https://arxiv.org/abs/2104.09864}

\bibitem{vaswani}
A.~Vaswani, N.~Shazeer, N.~Parmar, J.~Uszkoreit, L.~Jones, A.~N. Gomez, Ł.~Kaiser e I.~Polosukhin.
\newblock Attention is all you need.
\newblock En \emph{Advances in Neural Information Processing Systems 30 (NeurIPS)}, 2017.
\newblock \url{https://arxiv.org/abs/1706.03762}

\end{thebibliography}


\end{document}
